% Image credits for the photographs and AI illustrations of Book 3
% (University Chemistry, Year 2), Dutch edition. Machine-readable license
% records live in images/book3/CREDITS.md; keep this file in sync with
% frontmatter/image-credits-book3.tex.
\chapter*{Beeldverantwoording}

De tekeningen in dit boek zijn origineel. De foto's worden gebruikt onder
hun respectieve vrije licenties, met dank aan hun makers; de volledige
bronlinks en licentie-URL's van elke foto staan in
\texttt{images/book3/CREDITS.md} in de repository van het boek.
De gevarenpictogrammen zijn die van het Globally Harmonized System of
Classification and Labelling of Chemicals, zoals getekend voor het pakket
\texttt{ghsystem} van \TeX\ Live.

\bigskip
Naast de foto's en de originele tekeningen zijn sommige figuren
illustraties die met AI gegenereerd zijn, gemaakt met de beeldgeneratie
van OpenAI op basis van prompts die door de redactie van het boek
geschreven zijn, en vóór opname nagekeken op chemische juistheid. De
volledige prompt van elk gegenereerd beeld staat in
\texttt{images/book3/ai/PROMPTS.md} in de repository.

\bigskip
\noindent Foto's, per hoofdstuk:
\begin{itemize}\small
  \item Hfst.~1: Germain Henri Hess, portret door I.~I. Reimers (1837--1838),
        CC BY-SA 4.0; Marcellin Berthelot, onbekende auteur, publiek domein
        (beide Wikimedia Commons).
  \item Hfst.~2: Josiah Willard Gibbs, onbekende fotograaf, publiek domein
        (Wikimedia Commons).
  \item Hfst.~3: François-Marie Raoult, onbekende fotograaf, publiek domein
        (Wikimedia Commons).
  \item Hfst.~4: Jacobus Henricus van 't Hoff, door Nicola Perscheid (1904),
        publiek domein; Henry Le Chatelier, onbekende auteur, CC BY-SA 4.0;
        Fritz Haber en Carl Bosch, Nobel Foundation, publiek domein (alle
        Wikimedia Commons).
  \item Hfst.~6: hoogoven in Duisburg-Nord, door Dietmar Rabich, CC BY-SA
        4.0; thermietlassen van een rail, door PetrS., CC BY-SA 3.0 (beide
        Wikimedia Commons).
  \item Hfst.~7: destillatiekolom van een raffinaderij, door CEphoto, Uwe Aranas,
        CC BY-SA 3.0 (Wikimedia Commons).
  \item Hfst.~9: Michael Faraday, door Thomas Phillips (1842), publiek domein
        (Wikimedia Commons).
  \item Hfst.~10: Jaroslav Heyrovský, archief van het J.~Heyrovský-instituut
        voor fysische chemie, CC BY-SA 3.0 (Wikimedia Commons).
  \item Hfst.~11: een zuil van Volta, door GuidoB, CC BY-SA 3.0; Charles Martin Hall
        en Paul Héroult, onbekende auteurs, publiek domein (alle Wikimedia
        Commons).
  \item Hfst.~12: een koperen standbeeld, door Elcobbola, publiek domein; een opofferingsanode,
        door Zwergelstern, CC BY-SA 3.0 (beide Wikimedia Commons).
  \item Hfst.~13: Erwin Schrödinger, Nobel Foundation (1933), publiek domein
        (Wikimedia Commons).
  \item Hfst.~14: vloeibare zuurstof in een magneet, door Pieter Kuiper, publiek domein;
        Robert Mulliken en Friedrich Hund, door GFHund, CC BY 3.0 (beide Wikimedia
        Commons).
  \item Hfst.~16: August Kekulé (1873), onbekende auteur, publiek domein
        (Wikimedia Commons).
  \item Hfst.~18: oplossingen van overgangsmetalen, door Benjah-bmm27, publiek domein; Alfred
        Werner, Les Prix Nobel (1914), publiek domein (beide Wikimedia Commons).
  \item Hfst.~19: een robijnkristal, publiek domein; een smaragd, door Didier Descouens,
        CC BY-SA 3.0; kopersulfaatkristallen, door Stephanb, CC BY-SA 3.0 (alle
        Wikimedia Commons).
  \item Hfst.~20: Akira Suzuki, Ei-ichi Negishi en Richard Heck, door BloodIce, CC BY-SA
        4.0 (Wikimedia Commons).
  \item Hfst.~22: Charles Friedel (jaren 1890) en James Mason Crafts (\textit{Popular
        Science Monthly}, 1900), publiek domein (Wikimedia Commons).
  \item Hfst.~24: Michel Eugène Chevreul, door Nicolas Eustache Maurin, publiek domein
        (Wikimedia Commons).
  \item Hfst.~25: Alexander Borodin, onbekende auteur, publiek domein; Ludwig Claisen, door
        Veronica.villagrasa, CC BY-SA 4.0 (beide Wikimedia Commons).
  \item Hfst.~26: Robert Robinson, onbekende auteur, publiek domein (Wikimedia Commons).
  \item Hfst.~28: Elias James Corey, door Trvthchem, CC0 (Wikimedia Commons).
  \item Hfst.~29: Wallace Carothers in het laboratorium, onbekende fotograaf, publiek domein
        (Wikimedia Commons).
  \item Hfst.~30: Emil Fischer, Atelier Victoria, Berlijn, rond 1895, publiek domein (Wikimedia
        Commons).
  \item Hfst.~31: Michail Tsvet, onbekende auteur, publiek domein (Wikimedia Commons).
  \item Hfst.~32: Francis William Aston (1922), onbekende auteur, publiek domein; Gustav Kirchhoff,
        Robert Bunsen en Henry Roscoe (1862), publiek domein (beide Wikimedia Commons).
  \item Hfst.~34: William Sealy Gosset (1908), publiek domein (Wikimedia Commons).
%% next chapter here
\end{itemize}
\clearpage
